% Replacement text for Paper 63 (see debug/p63_pi_nonzero_memo.md for loci). Label kept to avoid breaking refs.
% (1) replaces l.1123-1131
\begin{proposition}[center derivatives reach the rung]
\label{prop:tree_attain_hom}
In the setting of Conjecture~\ref{conj:tree_attain}, let each $P_v$ be a
combination of $h_\ell(\partial)\nabla^{2j}$ with $\ell+2j\le2$, degrees mixed
or not (center derivatives only;\ scale derivatives are not covered).  If
$\Pi=\langle P_0((n-1)\omega)\prod_{v\neq0}P_v(-\omega)\rangle\neq0$
($\nabla^2\to1$), then $F_n$ lies in the $D$-module of the fibre, for
$3\le n\le8$;\ for $n\ge9$ (same setting otherwise) the two modules meet in a
nonzero submodule.  For homogeneous $P_v$,
$\Pi=(-(n-1))^{d_0}\langle\prod_v\sigma_v(-\omega)\rangle$.
\end{proposition}

% (2) replaces l.1153-1155 "its constant is C times ... nonzero by hypothesis."
its constant is $C\Pi$:\ every term of $P_v(\partial)K_0$, of any degree, has the
tail of $K_0$, so the leading angular factor is $P_v(-\omega)$, not its principal
part;\ Watson's lemma to second order on the two integrals of a channel-$L$ tail
gives $-s_0/k^2$ for every $L$, so that factor passes each vertex pointwise in
$\omega$;\ and $P_0(-\partial)$ acts on $e^{-(n-1)x}$ as $P_0((n-1)\omega)$.

% (3) replaces l.1169-1174 "For non-homogeneous P_v ... the s-type one)."
The hypothesis is $\Pi\neq0$, not the principal-symbol integral:\ a root
$1+\partial_z+\partial_z^2$ over a leaf $\partial_z+\tfrac12\nabla^2$ and a $p_z$
leaf has principal integral $0$ and $\Pi=4/5$, and is covered.

% (4) replaces the last sentence of conj:tree_attain (l.1119-1120)
When also $\Pi\neq0$ it is Proposition~\ref{prop:tree_attain_hom}, which covers
every homogeneous and every $s$-type tree;\ open is $\Pi=0$ (a leaf
$\partial_z^2-\tfrac13$ next to $s$ densities:\ $\Pi=0$, principal integral
$4\pi/3$).
